% ===================================================================
% DISCUSSION AND LIMITATIONS. ~0.5–1 page. Four honest bullets.
% Written from known open items in notes/open_questions.md and
% target.md §4. No overclaiming.
% ===================================================================

\paragraph{Cross-entropy extension.} Theorem~\ref{thm:general}
is proved in the quadratic-loss setting with $H_t$ a
positive-semidefinite rank-$r$ matrix. For LLM cross-entropy
loss, the natural $H_t$ is the Fisher information matrix, which
is rank-$r$ PSD near a LoRA fine-tuned point but not a projector.
Our proof uses only the per-task identity $H_tH_t^+H_t = H_t$
and the $O(d)$-equivariance of the hard distribution, both of
which survive verbatim under the Fisher-metric interpretation
(the $D_t$-freeness of Step 5 is the key). A fully rigorous
MSE-to-CE bridge via local quadratic approximation requires
controlling the second-order remainder across the $H_t$-ball;
we sketch this extension in the supplementary material and
leave a quantitative version for future work.

\paragraph{The shared-$V$ exponent gap.} In the floor-positive
regime (e.g.\ shared or substantially-overlapping $V_t$), the
explicit linear construction of \S\ref{sec:achv} achieves an
excess distortion decay of $2^{-2R/(\deff+|A|-1)}$, slower than
the lower bound's $2^{-2R/\deff}$. We verified numerically that
a random-codebook (non-linear) encoder achieves slopes strictly
steeper than the linear construction, so the truth lies
\emph{strictly between} the two bounds. Closing this gap — either
by sharpening the lower bound or by constructing a non-linear
encoder that matches the $2^{-2R/\deff}$ rate — is a concrete
open problem that we view as beyond the scope of this work.
Remark~\ref{rem:sharedv-gap} elaborates.

\paragraph{Scope of the hard distribution.} The lower bound is
established against $P^\star$ (Stiefel-random $V_t$, arbitrary
task-dependent spectra). Real LoRA fine-tunes exhibit structure
not captured by iid Stiefel random sampling: tasks that share
semantic similarity produce correlated subspaces; fine-tuning
induces anisotropic curvature; and the full-dimensional backbone
of the base model introduces rank-$d$ modes the bound ignores.
Extending the lower bound to structured task distributions, or
supplying a second construction that exploits task correlation
to beat our generic upper bound, are natural directions. Until
such extensions exist, the theorems of this paper should be read
as characterizing the \emph{information-theoretic} floor for
merging, not a tight prediction of any specific real-LLM merging
experiment.

\paragraph{Practical implications.} Two guidance points for
practitioners follow from the theorems. (i) The floor
$B^2(1 - \deff/(Tr))$ identifies when rate is wasted: if tasks
share substantial subspace overlap, no merging algorithm closes
the floor and additional bits buy little. This gives an
operational criterion for deciding when to merge at all versus
when to serve tasks from separate adapters. (ii) The
compression term $B^2 (\deff/(Tr)) 2^{-2R/\deff}$ is matched by
an explicit algorithm that runs in $O(\deff^2)$ time per merge
(the Gaussian-QR rotation dominates). For LoRA with $\deff = Tr$
in the thousands, this is negligible compared with the merging
overhead of existing heuristics; deployment cost is not a
blocker.
